% Generated draft: requires mathematical and source review.
\documentclass[11pt]{article}
\usepackage[margin=25mm]{geometry}
\usepackage{fontspec}
\setmainfont{Latin Modern Roman}
\usepackage{amsmath,amssymb,mathtools,unicode-math}
\setmathfont{Latin Modern Math}
\usepackage{enumitem,needspace,xurl,hyperref}
\setlength{\parindent}{0pt}
\setlength{\parskip}{5pt}
\setlength{\emergencystretch}{3em}
\newcommand{\N}{\mathbb{N}}
\newcommand{\Z}{\mathbb{Z}}
\newcommand{\Q}{\mathbb{Q}}
\newcommand{\R}{\mathbb{R}}
\newcommand{\C}{\mathbb{C}}
\newcommand{\F}{\mathbb{F}}
\newcommand{\A}{\mathbb{A}}
\newcommand{\PP}{\mathbb{P}}
\newcommand{\E}{\mathbb{E}}
\newcommand{\eps}{\varepsilon}
\newcommand{\dd}{\,\mathrm{d}}
\begin{document}
\section*{UnsolvedMath ID: 701744}
\textbf{Source problem code:} catalogue entry
\subsection{Definitions and mathematical statement}
Let \(G\) be a finite group and \(p\) a prime. Write \(\operatorname{Syl}_p(G)\) for the set of Sylow \(p\)-subgroups of \(G\), and suppose that
\[
s_p(G):=|\operatorname{Syl}_p(G)|=1+kp>1,
\qquad k\in\Z_{>0}.
\]
Choose \(P\in\operatorname{Syl}_p(G)\). Define
\[
G_p:=\{g\in G:\operatorname{ord}(g)=p^a\text{ for some }a\in\Z_{\ge0}\},
\]
so that the identity is included. The source attributes to Frobenius the assertion that
\[
f_p(G):=\frac{|G_p|}{|P|}\in\Z_{>0}.
\]
It further uses the integral parametrization
\[
f_p(G)=1+\ell(p-1),
\qquad \ell=\frac{f_p(G)-1}{p-1}.
\]
The intended question, attributed to P. Schmid, is whether every finite group \(G\) and prime \(p\) satisfying \(s_p(G)>1\) obey
\[
\ell^p\ge k^{p-1}.
\]
There is an unresolved inconsistency in the accompanying source assertions. The record states
\[
f_p(G)=1+\ell(p-1)\le k-\frac{k-1}{p-1},
\]
and says that usually \(\ell<k\), but that \(\ell=k\) if \(k<p\). In the latter case, the asserted equality gives \(f_p(G)=1+k(p-1)>k\), whereas the displayed upper bound is at most \(k\). These auxiliary claims are retained here as supplied, require source verification, and are not imposed as additional hypotheses restricting the universal question. The word ``usually'' has no specified mathematical quantifier.

For context, a finite group is called \(p\)-solvable if it has a subnormal series whose factors are \(p\)-groups or groups of order coprime to \(p\). The record reports that P. Gheri proved
\[
f_p(G)^p\ge s_p(G)^{p-1}
\]
when \(G\) is \(p\)-solvable. This is the supplied contextual result, with its original quantities and hypothesis; it is not identified with the requested inequality in \(\ell\) and \(k\).
\subsection{Short English statement}
For a finite group with more than one Sylow \(p\)-subgroup, write their number as \(1+kp\), and write the number of \(p\)-elements divided by the order of a Sylow \(p\)-subgroup as \(1+\ell(p-1)\). P. Schmid asks whether \(\ell^p\ge k^{p-1}\) always holds. An accompanying bound in the supplied record conflicts with its stated case \(\ell=k\) for \(k<p\), so that part of the formulation requires review.
\subsection{Sources}
\begin{itemize}
\item Supplied local UnsolvedMath registry record: permanent ID \(701744\), exact source code \texttt{catalogue entry}, title ``Kourovka Notebook Problem 20.86''. The registry labels the problem open; this is supplied source metadata, not an independently verified assessment.
\item \emph{The Kourovka Notebook (Unsolved Problems in Group Theory)}, Problem 20.86, proposed by P. Schmid in 2022. The record identifies the 20th issue (2022) and the 21st issue edition (2026), arXiv:1401.0300v48. The collection description names E. I. Khukhro and V. D. Mazurov as editors. These bibliographic identifications are supplied and have not been independently checked.
\item Frobenius: the divisibility assertion \(|P|\mid |G_p|\) is attributed to a theorem of Frobenius in the supplied statement. No specific bibliographic citation for that theorem is supplied.
\item P. Gheri, \emph{Ann. Mat. Pura Appl. (4)} \textbf{200} (2021), 1231--1243. The source reports the inequality \(f_p(G)^p\ge s_p(G)^{p-1}\) for \(p\)-solvable groups. No article title or link is supplied; this reported result has not been independently verified.
\end{itemize}
\end{document}
